% ============================================
% Supplementary Note 2: Assembly Effect Analysis
% ============================================
\subsection*{Supplementary Note 2: Assembly Effect Analysis}
\label{sec:suppl:note2}

To understand the mechanistic basis of community-level selection, we compared coalescence outcomes to a ``direct assembly'' control. In coalescence, two parental communities are first assembled separately from the species pool, then mixed; in direct assembly, all species from both pools are mixed simultaneously without prior assembly history.

During assembly, species compete and some go extinct. Those that survive to coexistence tend to have weaker mutual interactions on average. We quantified this assembly effect by measuring the mean pairwise interaction strength $\langle \alpha_{ij} \rangle$ among species at different stages (Supplementary Fig.~18). Across all tested interaction strengths ($\mu = 0$ to $1.2$), post-assembly communities exhibited substantially reduced mean interaction strength compared to the pre-assembly pool. This reduction becomes more pronounced at higher $\mu$, where stronger competition drives more species to extinction, leaving only those with weaker mutual interactions. The assembly process thus creates internal coherence within each parental community: species that survived together share compatible (less competitive) interactions.

We hypothesized that this shared history would increase Dominance probability compared to direct assembly. Simulations supported this prediction (Extended Data Fig.~7): across interaction strengths $\mu = 0.2$ to $1.0$, coalescence consistently produced higher Dominance fractions than direct assembly. At $\mu = 0.2$, coalescence yielded 37\% Dominance versus 12\% for direct assembly; at $\mu = 0.6$, 59\% versus 20\%; at $\mu = 1.0$, 65\% versus 22\%. Correspondingly, direct assembly showed higher Restructuring rates, as expected when species from different pools have not been pre-filtered for compatibility.

Experimental results corroborate these findings (Supplementary Fig.~19). Under Base and Nutr$+$ conditions, coalescence showed elevated Dominance rates compared to direct assembly, consistent with simulation predictions. The Nutr$-$ condition showed more variable results, likely due to weaker competitive interactions in nutrient-limited environments.

